\documentclass[10pt]{article}

\usepackage[margin=1in]{geometry}
\usepackage{booktabs}
\usepackage{caption}
\usepackage{graphicx}
\usepackage{hyperref}
\usepackage{microtype}
\usepackage{tabularx}

\graphicspath{{../../../../results/paper/qualitative-figure-panels/}}

\title{FocusParse: Evidence Packets for High-Resolution Document Question Answering}
\author{Anonymous submission}
\date{}

\begin{document}
\maketitle

\begin{abstract}
High-resolution technical and financial PDFs expose a persistent failure mode
for multimodal document question answering: whole pages preserve context but
compress away small labels, footnotes, and chart/table details, while tight
crops recover readability but detach evidence from the local context needed to
interpret it. We study this failure mode using parser-bench, a
citation-required visual QA benchmark over finance reports and technical
datasheets where answers are evaluated with both correctness and page/box
grounding. We introduce FocusParse, a structured harness that routes to likely
pages, localizes candidate regions, inspects high-resolution crops, expands
selected evidence with role-labeled neighboring context, and answers only from
compact evidence packets. In a matched seven-method run on the pinned 148-row
paper slice, FocusParse +4 reaches 61.5\% overall accuracy, 66.3\% datasheet
accuracy, and 51.1\% finance accuracy, compared with 43.9\% for a base VLM and
6.1--18.2\% for generic tool-agent comparators. The same run shows high page
recall (0.914), strong region overlap (0.857 mean BBox IoU), and low lazy
answering (3.4\%). The resulting thesis is that performance on dense finance
and datasheet documents depends not merely on bigger context windows or more
tools, but on constructing the right evidence unit before reasoning: a readable
region plus the captions, legends, headers, footnotes, axes, and cross-page
references that make the region meaningful.
\end{abstract}

\section{Introduction}

Document AI systems increasingly need to answer questions over high-resolution
PDFs rather than simply transcribe them. In finance reports and technical
datasheets, the answer often depends on small plotted labels, multi-panel
figures, footnotes, min/typ/max tables, timing diagrams, and symbol dictionaries
spread across multiple pages. This creates a resolution-context tradeoff. A
full-page VLM call can preserve surrounding layout but may make fine details
unreadable. A tight crop can make the relevant text or visual element readable
but may lose the caption, legend, table header, or footnote that determines what
the crop means.

FocusParse is built around the claim that the useful intermediate object for
this setting is an \emph{evidence packet}. An evidence packet binds a localized
region to its supporting context: local crop, text/OCR snippets, role-labeled
neighbor crops, provenance, confidence, and optional extracted structure. The
method is deliberately narrower than a general-purpose parser. It is a harness
for query-conditioned evidence construction in dense documents where the answer
must be tied to supporting pages and boxes.

This paper makes four contributions. First, it frames parser-bench as a
region-grounded QA setting for concentrated finance and datasheet documents.
Second, it describes FocusParse, an evidence-localization-first harness with
explicit inspect and expand stages. Third, it compares base VLMs, generic ReAct
loops, generic tool agents, and FocusParse under matched tools and scorers.
Fourth, it analyzes answer accuracy, cost per correct, page recall, BBox IoU,
lazy-answer rate, and qualitative packet traces.

\section{Related Work}

Document parsing benchmarks show that document pages contain structured visual
regions and that parsing quality should be evaluated beyond plain OCR
\cite{doclaynet2022,pubtabnet2020,gtefintabnet2021,omnidocbench2025,mpdocbenchparse2026,parsebench2026}.
MPDocBench-\allowbreak Parse is useful because it moves parsing evaluation toward practical
multi-page documents, semantic continuity, hierarchy recovery, and visual
content preservation. External ParseBench is especially relevant because it
evaluates agent-facing parsing along tables, charts, content faithfulness,
semantic formatting, and visual grounding. FocusParse asks a different but
adjacent question: not whether a page or multi-page document can be fully
parsed, but whether a QA system can select the answer-supporting evidence for a
query, attach the local context needed to interpret it, and cite the supporting
page/box. Throughout this draft, \texttt{ParseBench} refers to the RunLlama
parsing benchmark, while \texttt{parser-bench} refers to the Gabriel benchmark
used for FocusParse evaluation.

Document QA and financial QA benchmarks establish the need for multimodal and
domain-specific reasoning. DocVQA, InfographicVQA, and MMLongBench-Doc study
questions over document images and long multimodal contexts
\cite{docvqa2021,infographicvqa2021,mmlongbenchdoc2024}. FinQA, TAT-QA,
FinanceBench, and FinRAGBench-V emphasize numerical reasoning, open-book
financial QA, and visual citations
\cite{finqa2021,tatqa2021,financebench2023,finragbenchv2025}. Parser-bench
differs by making high-resolution region localization, multi-piece context
collection, and page/crop citation part of the QA contract.

Chart and table QA benchmarks such as ChartQA and PlotQA show that charts
require both visual perception and reasoning \cite{chartqa2022,plotqa2020}.
FocusParse places this challenge inside long PDFs where chart interpretation
often depends on neighboring captions, axes, legends, footnotes, or table
references. The chart-reading problem is therefore inseparable from evidence
collection: the model must read the mark and recover the nearby semantics that
make the mark meaningful.

The closest methodological neighbors are AgenticOCR and DocLens. AgenticOCR
argues that page-level chunking introduces extraneous context and dilutes
salient evidence, reframing OCR as query-conditioned on-demand extraction
\cite{agenticocr2026}. DocLens identifies evidence localization as a central
failure mode in long visual document understanding and navigates from documents
to pages and elements before answer adjudication \cite{doclens2025}. FocusParse
shares the active evidence-acquisition motivation, but studies a controlled
stage machine with explicit typed evidence packets and role-labeled context
expansion. The contrast is the object handed to the reasoner: a compact packet
containing a readable crop plus captions, legends, headers, footnotes, axes, or
neighboring regions, rather than a loose transcript of tool observations.

Generic ReAct-style agents provide a natural baseline for tool use
\cite{react2023}. The key comparison is not tools versus no tools. It is
unconstrained tool use versus structured evidence construction. Generic loops
may inspect documents, but in dense PDFs they can also terminate lazily, repeat
irrelevant calls, or answer from noisy transcripts that contain weak citations.

\section{Parser-Bench Dataset}

Parser-bench is a citation-required visual QA benchmark over full-resolution
finance reports and technical datasheets. Each example contains a question, a
gold answer, answer-type metadata, source document metadata, rendered page
images, supporting pages, supporting bounding boxes, optional alternate boxes,
question-family labels, difficulty axes, and reasoning-chain metadata. The
current paper slice is the canonical post-filter validation materialization:
148 rows, 147 unique IDs, 44 source PDFs, 101 datasheet examples, and 47
finance examples.

\begin{table}[t]
\centering
\caption{Canonical parser-bench paper slice used for the final matched run.}
\label{tab:dataset}
\begin{tabular}{lrrrrr}
\toprule
Slice & Rows & Unique IDs & PDFs & Datasheet & Finance \\
\midrule
Canonical paper slice & 148 & 147 & 44 & 101 & 47 \\
\bottomrule
\end{tabular}
\end{table}

The slice is small but deliberately concentrated. It contains 70 multi-region
examples, 34 multi-page examples, 112 visually grounded examples, and 49 rows
that are both visual and multi-region. The most frequent question families are
axis-value interpolation, near-miss distractors, cross-page continuation,
min/typ/max disambiguation, and chart/table cross-reference. These are exactly
the families where inspection and context expansion matter: visual labels must
be read at high resolution, then connected to surrounding document semantics.

The current pinned materialization records a fixed Hugging Face revision,
dataset fingerprint, and benchmark JSONL SHA-256 in the artifact manifest. The
full identifiers are reported with the released evaluation package rather than
treated as moving URLs.

\section{Method: FocusParse}

FocusParse is an eight-stage workflow:
\begin{center}
\texttt{PLAN $\rightarrow$ ROUTE\_PAGES $\rightarrow$ LOCALIZE $\rightarrow$ RERANK $\rightarrow$ INSPECT $\rightarrow$ EXPAND\_CONTEXT $\rightarrow$ ANSWER $\rightarrow$ VERIFY}.
\end{center}
The stages communicate through typed events rather than arbitrary scratchpad
state. Planning predicts question family and evidence types. Routing chooses
candidate pages. Localization and reranking propose visual/textual regions.
Inspection turns selected regions into readable evidence by cropping,
extracting text/OCR, and optionally using region-specific tools. Expansion
attaches bounded context such as captions, legends, headers, axes, footnotes,
and nearby table text. Answering receives evidence packets rather than raw
pages. Verification checks whether the answer is supported, whether citations
are grounded, and whether a bounded repair action is needed.

The central contract is \texttt{EvidencePacket}. A packet contains page,
normalized box, region type, provenance, crop references, text-layer and OCR
snippets, linked neighbor crop references, neighbor-role labels, optional
chart/table extraction, confidence, and selection metadata. This design creates
an explicit bottleneck: the reasoner can only use compact evidence assembled by
the upstream stages. That bottleneck is the mechanism the paper tests.

\section{Experiments}

The final matched comparison has seven rows: a base VLM, ReAct +2 tools, ReAct
+4 tools, generic agent +2 tools, generic agent +4 tools, FocusParse +2 tools,
and FocusParse +4 tools. The +2 tool set exposes \texttt{inspect\_region} and
\texttt{get\_text\_layer}; the +4 tool set adds \texttt{expand\_context} and
\texttt{run\_python}. Metrics include answer accuracy with 95\% bootstrap
confidence intervals, cost per correct answer, mean latency, page recall, mean
BBox IoU, lazy-answer rate, and qualitative packet traces.

\section{Results}

\begin{table*}[t]
\centering
\caption{Final matched seven-method result on the pinned 148-row parser-bench slice.}
\label{tab:main-results}
\small
\begin{tabular}{lrrrrrr}
\toprule
Method & Overall & Datasheet & Finance & \$/correct & Page recall & BBox IoU \\
\midrule
Base VLM & 43.9 & 47.5 & 36.2 & 0.0102 & 0.848 & 0.000 \\
ReAct +2 & 18.2 & 22.8 & 8.5 & 0.1341 & 0.819 & 0.336 \\
ReAct +4 & 16.2 & 20.8 & 6.4 & 0.1946 & 0.727 & 0.309 \\
Agent baseline +2 & 8.1 & 10.9 & 2.1 & 0.1324 & 0.000 & 0.000 \\
Agent baseline +4 & 6.1 & 7.9 & 2.1 & 0.1811 & 0.000 & 0.000 \\
FocusParse +2 & 60.1 & 64.4 & 51.1 & 0.0282 & 0.929 & 0.897 \\
FocusParse +4 & 61.5 & 66.3 & 51.1 & 0.0248 & 0.914 & 0.857 \\
\bottomrule
\end{tabular}
\end{table*}

FocusParse +4 improves over the base VLM by 17.6 overall points, 18.8
datasheet points, and 14.9 finance points. It improves over generic ReAct and
generic agent baselines by 43--55 overall points while remaining much cheaper
per correct than those tool-agent comparators. Base VLM remains cheapest per
correct, so the claim is an accuracy and localization gain, not a lowest-cost
claim.

Diagnostics show that FocusParse +4 builds 1184 evidence packets, calls
\texttt{expand\_context} on every example, attaches 9.72 neighbors on average,
and cites packets with 100\% text coverage and 84.4\% linked-context coverage.
Generic ReAct calls tools but does not construct evidence packets, while the
generic agent baselines mostly terminate without useful tool use.

\subsection{Mechanism Ablations}

Table~\ref{tab:ablations} summarizes the current n=148 mechanism ablations.
The strongest positive signal is query-conditioned reranking: restoring rerank
improves accuracy from 56.1\% to 61.5\% and substantially improves grounding.
Restoring \texttt{expand\_context} gives a smaller +2.0 point gain. In contrast,
restoring verifier-directed repair moves 62.8\% to 61.5\% in this matched run,
and restoring answer-shape repair moves 64.2\% to 61.5\%, so the repair
mechanisms are best treated as mixed or negative controls rather than sources of
the paper gain.

\begin{table}[t]
\centering
\caption{Mechanism ablations against the final FocusParse +4 row.}
\label{tab:ablations}
\small
\begin{tabular}{lrrrrr}
\toprule
Condition pair & Baseline & Full +4 & $\Delta$ & Rec. & Reg. \\
\midrule
No rerank $\rightarrow$ rerank & 56.1 & 61.5 & +5.4 & 19 & 11 \\
No \texttt{expand\_context} $\rightarrow$ expand & 59.5 & 61.5 & +2.0 & 14 & 11 \\
Repair off $\rightarrow$ repair on & 62.8 & 61.5 & -1.4 & 9 & 11 \\
Answer-shape repair off $\rightarrow$ on & 64.2 & 61.5 & -2.7 & 9 & 13 \\
\bottomrule
\end{tabular}
\end{table}

\subsection{Failure Taxonomy}

The final FocusParse +4 row is correct on 91/148 examples and incorrect on
57/148 examples. Table~\ref{tab:failures} summarizes the primary error buckets
computed from the final \texttt{per\_example.jsonl}.

\begin{table}[t]
\centering
\caption{Primary failure buckets for the final FocusParse +4 run.}
\label{tab:failures}
\small
\begin{tabular}{lrrrr}
\toprule
Failure bucket & Count & Datasheet & Finance & Error share \\
\midrule
Verifier unsupported & 26 & 16 & 10 & 45.6 \\
Partial localization & 12 & 6 & 6 & 21.1 \\
Localization miss & 8 & 5 & 3 & 14.0 \\
Reasoning/extraction & 7 & 4 & 3 & 12.3 \\
Lazy/no bbox & 4 & 3 & 1 & 7.0 \\
\bottomrule
\end{tabular}
\end{table}

The dominant remaining error is not generic tool laziness: only 4/57 errors are
lazy/no-bbox failures. The largest bucket is verifier-unsupported answers after
evidence collection, followed by incomplete or missed localization. This
suggests that the next lift is likely better verification, evidence
adjudication, and multi-region localization rather than simply adding tools.

\section{Qualitative Analysis}

\begin{figure*}[t]
\centering
\includegraphics[width=\textwidth]{figure3-finance-latvia-evidence-binding.png}
\caption{Finance evidence binding. The answer requires joining two chart-panel
conditions to a country-code table. Same-revision baselines expose the failure
mode: Base VLM returns unanswerable, generic agents often choose Estonia or
emit raw tool transcripts, while FocusParse +4 answers \emph{Latvia} with all
three cited pages and near-perfect region overlap.}
\label{fig:latvia}
\end{figure*}

\begin{figure*}[t]
\centering
\includegraphics[width=\textwidth]{figure4-datasheet-jesd204b-evidence-binding.png}
\caption{Datasheet evidence compaction. Several baselines can state or nearly
state the integer answer, but they fail to preserve complete page/box grounding.
FocusParse +4 answers \emph{6} with the timing diagram and cross-page
corroborating block diagrams retained as citable evidence packets.}
\label{fig:jesd}
\end{figure*}

The qualitative examples make the main metric story inspectable. In the Latvia
finance example, the answer turns on linking chart labels from pages 108 and
109 to the \texttt{LV -> Latvia} country-code table on page 8. FocusParse +4
has page recall 1.0, BBox IoU 0.99999, and three citations. ReAct +2 recovers
the answer but with weak region grounding, while generic agents either choose
Estonia or emit invalid raw transcripts. In the JESD204B datasheet example,
baselines can sometimes state the answer \emph{6}, but FocusParse is the row
that preserves complete page recall and near-perfect region overlap across the
timing and block-diagram evidence.

\section{Limitations}

The current dataset slice is small and concentrated: 148 rows with 47 finance
examples. The paper therefore reports confidence intervals and avoids broad
claims about all document QA. A stronger 66.9\% FocusParse checkpoint is
documented in research notes, but the raw run directory is not present in the
current checkout; it is not used as a headline claim here. The qualitative
panels and failure-taxonomy table are draft analysis artifacts; final styling
depends on the selected venue template and page budget. Zotero export is
currently blocked because Zotero Desktop/local API is unavailable on this
profile. Finally, FocusParse is a research harness, not a general document
parser or production service.

\section{Conclusion}

FocusParse studies a concrete mechanism for high-resolution document QA:
construct the evidence before asking the model to answer. On the pinned 148-row
parser-bench paper slice, FocusParse +4 reaches 61.5\% accuracy with strong
page and region grounding, outperforming both a no-tool base VLM and generic
tool-agent comparators. The more important mechanism signal is that structured
inspect and expand stages produce compact packets that preserve both
readability and interpretive context. This supports the paper's central claim:
for dense finance and datasheet documents, better harnesses are not just models
with larger context windows or more tools. They are systems that build the
right evidence object before reasoning.

\bibliographystyle{plain}
\bibliography{references}

\appendix
\section{Artifact Provenance}

The main result table, composed qualitative panels, and same-revision
qualitative baseline predictions are sourced from the May 24 result package
recorded in the public experiment runbook. A slim review archive is available at
\texttt{results/paper/submission-review-package/}; the exact archive path is
recorded in the package manifest. All claims in the abstract use that final
matched package.

\end{document}
